\section{Conclusion}

Paper 14 now establishes a native finite theorem of Thalean relational
saturation.

The result does not identify native loading with iteration of the
Thalean quadratic.

Instead, it separates three exact finite operations:
\[
\text{selection},
\qquad
\text{retention},
\qquad
\text{concentration}.
\]

Native selection is supplied by the certified registered carrier.

For a live prefix \(p\), the lawful completion family contracts
monotonically:
\[
8\longrightarrow2\longrightarrow1.
\]

Registered-future saturation occurs when
\[
|\mathcal C(p)|=1.
\]

Future uniqueness does not erase complete relational distinction.

At the two-edge threshold there remain \(120\) distinct contexts with
\(120\) distinct registered completions, and the terminal completion
register retains two native holonomy designations over each closing
chamber.

Transport also creates the context in which native designation becomes
operational.

If \(\Delta\) is the induced context-activated designation response
and \(H\) the ten-dimensional five-address-hidden projector, then
\[
\Delta^{\mathsf T}\Delta
=
\Delta\Delta^{\mathsf T}
=
\frac34H,
\]
with
\[
\rho\Delta=0.
\]

Thus the native response is exactly supported on a retained relational
sector that is invisible to the declared coarse readout.

Native loading does not itself perform common-mode contrast
concentration.

For the canonical completion-occupancy register,
\[
\zeta:
\frac{31}{55}
\longrightarrow
\frac{41}{65}
\longrightarrow
\frac{7}{10},
\]
and
\[
\rho^2:
\frac{31}{24}
\longrightarrow
\frac{41}{24}
\longrightarrow
\frac73.
\]

Every one of the \(180\) strict native loading extensions increases
these contrast statistics.

Native loading therefore sharpens the selected future.

The normalized Thalean quadratic supplies the distinct concentrating
operation:
\[
\widehat Q=\frac1{98}Q,
\qquad
Q=MM^{\mathsf T}=A^3+2A^2+2I.
\]

Its common mode is fixed, every centered mode contracts, and
\[
R(\widehat Q^{\,k}x)
\le
\left(\frac9{49}\right)^kR(x).
\]

Every finite iterate remains invertible, even though
\[
\widehat Q^{\,k}\longrightarrow P_0.
\]

The native loading response and the centered quadratic belong to the
same certified transport-response algebra, but they are not the same
operator and do not perform the same finite action.

\subsection{What is theorem}

The following statements are established internally:

\begin{enumerate}[label=(\roman*)]

\item native relational loading is monotone contraction of the lawful
completion family;

\item the exact loading profile is
\[
8\to2\to1;
\]

\item registered-future saturation occurs at the two-edge threshold;

\item unique future does not imply unique complete context;

\item context activates native \(b/ab\) designation after one admitted
carrier step;

\item the induced designation response has rank ten;

\item
\[
\Delta^{\mathsf T}\Delta
=
\Delta\Delta^{\mathsf T}
=
\frac34H;
\]

\item the retained response is invisible to the five-address readout;

\item the canonical completion register strictly increases Paper 14
contrast along all \(180\) strict loading extensions;

\item native future saturation and contrast saturation are distinct;

\item the normalized Thalean quadratic contracts centered contrast by
the sharp factor \(9/49\);

\item every finite quadratic iterate remains invertible;

\item the native loading response and the centered quadratic inhabit
the same exact transport-response algebra.

\end{enumerate}

Together these statements close the native finite
relational-saturation theorem.

\subsection{What remains conditional internally}

The following are not promoted to theorem:

\begin{itemize}

\item a unique clean-release optimum;

\item a universal overcompression threshold;

\item a dynamical cycle specifying when native selection is followed
by or coupled to quadratic concentration;

\item a globally closed inter-face sigma;

\item a universal finite information-capacity law.

\end{itemize}

The open sigma problem belongs to a stronger global inter-face
transport construction.

It is not a missing premise of the finite saturation theorem.

\subsection{The theorem ladder}

The closed finite chain is now

\[
\text{native registered carrier}
\]
\[
\downarrow
\]
\[
\text{context activation}
\]
\[
\downarrow
\]
\[
8\to2\to1
\]
\[
\downarrow
\]
\[
\text{registered-future saturation}
\]
\[
\downarrow
\]
\[
\Delta^{\mathsf T}\Delta=\frac34H
\]
\[
\downarrow
\]
\[
\text{retained projection-relative distinction}.
\]

Separately,
\[
\text{Thalean incidence algebra}
\]
\[
\downarrow
\]
\[
\widehat Q
\]
\[
\downarrow
\]
\[
\text{centered contrast contraction}
\]
\[
\downarrow
\]
\[
\text{finite-stage non-erasure}.
\]

The two chains meet in the common certified centered
transport-response algebra.

They are not collapsed into one operation.

\subsection{Physical boundary}

Nothing in the finite theorem identifies relational saturation with:

\begin{itemize}

\item physical gravity;

\item Newtonian or relativistic gravitational dynamics;

\item black-hole collapse;

\item an event horizon;

\item singularity avoidance;

\item black-hole entropy;

\item Hawking radiation;

\item Planck-scale information density;

\item physical proper time;

\item the cosmic microwave background;

\item or a Standard Model interaction.

\end{itemize}

Those require separate physical bridge theorems and empirical tests.

\subsection{Final statement}

The central finite lesson is:

\begin{quote}
Relational saturation is reached when admissible future freedom
becomes unique while complete relational distinction remains retained
behind registration.
\end{quote}

The exact native signature is
\[
|\mathcal C|:
8\longrightarrow2\longrightarrow1,
\]
together with
\[
\Delta^{\mathsf T}\Delta=\frac34H,
\qquad
\rho\Delta=0.
\]

The exact quadratic signature is
\[
R(\widehat Q^{\,k}x)
\le
\left(\frac9{49}\right)^kR(x)
\]
with finite-stage invertibility.

The mature theory therefore does not say that greater native loading
automatically produces lower contrast.

It says something more structured:

\begin{quote}
Selection sharpens a lawful future.
Context retains distinctions hidden by coarse registration.
Quadratic aggregation can independently suppress registered contrast
without erasing the retained state.
\end{quote}

Whether nature realizes any physical analogue of this finite
architecture remains an open question.
